\begingroup\small
\begin{longtable}{@{}p{0.44\textwidth}p{0.52\textwidth}@{}}
\toprule Declaration & Exact scope\\\midrule
\endhead
\multicolumn{2}{@{}l}{\textbf{\decl{H12PrincipalDivisor}}}\\[4pt]
\decl{principal5_cube} & The lattice endomorphism cubed is multiplication by seven.\\
\decl{lift_identity} & Exact lifted image with both residue carries.\\
\decl{digit_is_remainder} & Digit equals the Euclidean remainder modulo seven.\\
\decl{digit_range} & Digit lies in the interval from zero to six.\\
\decl{unique_lift_digit} & Uniqueness under the residue range and lift equation.\\
\decl{lift_preserves_fifth_weight} & The lifted fifth-residue weight is unchanged.\\
\decl{wall_bounds} & Wall lies between minus one and one.\\
\decl{wall_values} & Wall is minus one, zero, or one.\\
\decl{order_difference} & Coefficient is the difference of lifted/base cone weights.\\
\decl{order_values} & Coefficient is one of the five possible values.\\
\decl{tau_ray_order} & Positive tau-ray coefficients equal four.\\
\decl{s_ray_order} & Nonzero s-ray coefficients equal minus four.\\
\decl{negative_integer_order} & Positive-index negative-integer coefficients equal four.\\
\decl{zero_order} & Coefficient at the origin equals zero.\\
\addlinespace[4pt]
\multicolumn{2}{@{}l}{\textbf{\decl{H12CubicDivisorGrowth}}}\\[4pt]
\decl{positive_pivot} & The strict chamber inequality makes the pivot positive.\\
\decl{chamber_wall} & The specified sign chamber has wall one.\\
\decl{chamber_order} & Chamber coefficient equals its fifth-residue weight.\\
\decl{zero_family_order} & Positive-family coefficient equals four.\\
\decl{pole_family_order} & Pole-family coefficient equals minus one.\\
\decl{zero_family_injective} & Equal zero points recover the three free indices.\\
\decl{pole_family_injective} & Equal pole points recover all four indices.\\
\decl{families_disjoint} & No zero point equals a pole point with residue one to four.\\
\decl{zero_box_bounds} & Closed upper coordinate bound for the zero box.\\
\decl{pole_box_bounds} & Strict upper coordinate bound for the pole box.\\
\decl{origin_not_in_families} & Positive first index excludes the origin in both families.\\
\addlinespace[4pt]
\bottomrule\end{longtable}
\endgroup
